% =====================================================================
\parte{Parte B. Problemas}{cap:problemas}
% =====================================================================

Seis problemas de desarrollo, uno por módulo (70 puntos en total). La puntuación de cada apartado figura entre corchetes. Tiempo orientativo: 90 minutos.

% ---------------------------------------------------------------------
\begin{ejercicio}{2}{eje:p1}
\textbf{Modelo de bigramas (M01). 12 puntos.}

El corpus de entrenamiento está formado por las palabras «mar», «ama» y «rama». El vocabulario es $\V=\{\fin,\texttt{a},\texttt{m},\texttt{r}\}$, donde \fin\ marca el inicio y el fin de cada palabra.
\begin{enumerate}[label=(\alph*)]
  \item \pts{3} Construya la matriz de conteos $N$ y la matriz de probabilidades de máxima verosimilitud $P$. ¿Cuántos bigramas hay en total?
  \item \pts{3} Calcule $P(\text{mar})$, $P(\text{rama})$ y $P(\text{mama})$. Explique por qué «mama», que no aparece en el corpus, recibe la misma probabilidad que «ama».
  \item \pts{3} Calcule la pérdida media (entropía cruzada) del modelo sobre el corpus de entrenamiento, en nats y en bits, y su perplejidad. Compárela con la del modelo uniforme sobre $\V$ y valore la comparación.
  \item \pts{2} Aplique suavizado de Laplace ($\alpha=1$) y calcule $P(\texttt{a}\mid\texttt{m})$ y $P(\texttt{r}\mid\texttt{r})$.
  \item \pts{1} Con el orden $(\fin,\texttt{a},\texttt{m},\texttt{r})$ y el método de la inversa de la función de distribución, determine el token que sigue a «\texttt{a}» para $u=0{,}35$, $u=0{,}55$ y $u=0{,}9$.
\end{enumerate}
\fuente{elaboración propia, con la notación de Jurafsky y Martin~\cite[cap.~3]{jm3}.}
\end{ejercicio}

% ---------------------------------------------------------------------
\begin{ejercicio}{2}{eje:p2}
\textbf{Embeddings y analogías (M02). 10 puntos.}

Se consideran los siguientes embeddings de juguete en $\R^{3}$; $e_w$ denota el embedding de la palabra $w$.
\begin{center}
\begin{tabular}{@{}lc@{\qquad}lc@{\qquad}lc@{}}
\toprule
Palabra & Vector & Palabra & Vector & Palabra & Vector \\
\midrule
rey & $(3,2,1)$ & mujer & $(1,-2,2)$ & princesa & $(2,-1,2)$ \\
hombre & $(0,2,2)$ & reina & $(3,-2,1)$ & niña & $(0,-1,2)$ \\
\bottomrule
\end{tabular}
\end{center}
\begin{enumerate}[label=(\alph*)]
  \item \pts{2} Calcule $e_{\text{rey}}\cdot e_{\text{reina}}$, $\norm{e_{\text{rey}}}$ y $\cos(e_{\text{rey}},e_{\text{reina}})$.
  \item \pts{4} Calcule $q=e_{\text{rey}}-e_{\text{hombre}}+e_{\text{mujer}}$ y su coseno con $e_{\text{reina}}$, $e_{\text{princesa}}$ y $e_{\text{niña}}$. ¿Qué palabra devuelve la analogía? ¿Qué coordenada parece codificar el género?
  \item \pts{2} Aplique $W=\begin{pmatrix}1&0&0\\0&0&1\end{pmatrix}$ a $e_{\text{rey}}$ y a $e_{\text{reina}}$. ¿Qué coseno tienen las imágenes? Determine el núcleo de $W$ y explique qué información destruye.
  \item \pts{2} \emph{Estimación.} La capa de salida de un LLM multiplica un vector de dimensión $1024$ por una matriz de $50\,000\times 1024$ para obtener los logits. ¿Cuántas operaciones de coma flotante (productos y sumas) cuesta por token? Con un procesador que ejecuta $10^{12}$ operaciones útiles por segundo, ¿cuánto tarda?
\end{enumerate}
\fuente{elaboración propia, a partir del ejemplo de analogías de Mikolov \emph{et al.}~\cite{mikolov-b}.}
\end{ejercicio}

% ---------------------------------------------------------------------
\begin{ejercicio}{2}{eje:p3}
\textbf{Autovalores, SVD y bajo rango (M03). 12 puntos.}

Sean $A=\begin{pmatrix}3&1\\1&3\end{pmatrix}$ y $B=\begin{pmatrix}0&-2\\3&0\end{pmatrix}$.
\begin{enumerate}[label=(\alph*)]
  \item \pts{3} Calcule los autovalores de $A$ y una base ortonormal de autovectores. ¿Son ortogonales los autovectores? Justifique por qué.
  \item \pts{2} Calcule $A^{-1}$, compruebe que $AA^{-1}=I$ y relacione los autovalores de $A^{-1}$ con los de $A$.
  \item \pts{3} Calcule la SVD $B=U\Sigma V\T$ en dos pasos: primero $V$ y $\Sigma$ a partir de $B\T B$; después las columnas de $U$ mediante $u_i=Bv_i/\sigma_i$. Obtenga la mejor aproximación de rango~1, $B_1$, el error $\norm{B-B_1}_F$ y la fracción de energía $\sigma_1^{2}/(\sigma_1^{2}+\sigma_2^{2})$ que retiene. ¿Por qué la SVD de $B$ no puede leerse de sus autovalores, como en el caso de $A$?
  \item \pts{4} \emph{Estimación.} Una matriz de pesos de $2048\times 1024$ se almacena como aproximación de rango $k=64$ (factores $U_k$, $V_k$ y los $k$ valores singulares). ¿Cuántos números se almacenan y qué fracción del original suponen? ¿A partir de qué $k$ deja de compensar? Compárelo con LoRA de rango $r=8$ sobre una matriz de $4096\times 4096$: ¿cuántos parámetros se entrenan y qué porcentaje de la matriz representan?
\end{enumerate}
\fuente{apartado (c) adaptado de MIT 18.06, primavera de 2010, examen~3, problema~3(c)~\cite{mit1806ex3} (SVD en dos pasos a partir de $A\T A$, con otra matriz); apartado (d), del ejercicio de compresión y LoRA de los apuntes del M03 y de Hu \emph{et al.}~\cite{hu-lora}.}
\end{ejercicio}

% ---------------------------------------------------------------------
\begin{ejercicio}{3}{eje:p4}
\textbf{Gradiente de la softmax con entropía cruzada (M04). 14 puntos.}

Una capa lineal sin sesgo produce los logits $z=Wx$, con
\[
W=\begin{pmatrix}1&0\\0&1\\1&-1\end{pmatrix},\qquad x=\begin{pmatrix}1\\2\end{pmatrix},
\]
y el token correcto es el primero ($c=1$). Sean $p=\softmax(z)$, $y$ el vector one-hot de $c$ y $L=-\ln p_c$.
\begin{enumerate}[label=(\alph*)]
  \item \pts{4} \emph{Derivación.} Demuestre que $L$ coincide con la entropía cruzada $H(y,p)=-\sum_k y_k\ln p_k$ y que puede escribirse como $L=-z_c+\ln\sum_j e^{z_j}$. A partir de esta expresión, demuestre que $\partial L/\partial z_k=p_k-y_k$.
  \item \pts{3} Calcule $z$, $p$ y $L$.
  \item \pts{4} Calcule $\nabla_z L$, $\partial L/\partial W$ y $\partial L/\partial x$ mediante la regla de la cadena, indicando la forma de cada resultado.
  \item \pts{2} Dé un paso de descenso de gradiente sobre $W$ con tasa de aprendizaje $\eta=0{,}5$ y calcule la nueva pérdida.
  \item \pts{1} Compruebe $\partial L/\partial W_{11}$ con una diferencia centrada de paso $h=0{,}01$ y calcule el error relativo $r(a,n)=\abs{a-n}/\max(\abs{a},\abs{n})$.
\end{enumerate}
\fuente{apartados (a) y (c) adaptados de Stanford CS224N, tarea~2, preguntas~1(a), 1(b) y 1(d)--(e)~\cite{cs224n-a2} (equivalencia entre la pérdida \emph{naive-softmax} y la entropía cruzada, y gradientes respecto de los vectores y de la matriz de salida); el resto, de elaboración propia.}
\end{ejercicio}

% ---------------------------------------------------------------------
\begin{ejercicio}{2}{eje:p5}
\textbf{Estimación y regla de Bayes (M05). 11 puntos.}
\begin{enumerate}[label=(\alph*)]
  \item \pts{4} \emph{Derivación.} Un modelo muy simple supone que, tras cada letra, el nombre termina con probabilidad constante $q$, de modo que la longitud $L$ (en letras) sigue la distribución geométrica $P(L=k)=(1-q)^{k-1}q$, $k\ge 1$. Se observan las longitudes $3$, $4$, $5$, $5$, $6$ y $7$. Obtenga el estimador de máxima verosimilitud $\hat q$, compruebe que es un máximo y calcule su valor, $P(L=3)$ y $\E[L]=1/\hat q$.
  \item \pts{3} Un cajón contiene tres dados de 4 caras, uno de 6 y uno de 12, todos equilibrados. Se extrae un dado al azar, se lanza una vez y sale un~5. Construya la tabla de Bayes (a priori, verosimilitud, producto y a posteriori) para el tipo de dado. Calcule la probabilidad predictiva a priori de obtener un~5 y la probabilidad de que un segundo lanzamiento del mismo dado dé un~6.
  \item \pts{4} \emph{Estimación numérica.} Un modelo acierta el $78\,\%$ de un banco de pruebas de $600$ preguntas. Calcule el error típico de la exactitud y el intervalo de confianza aproximado del $95\,\%$. ¿Cuántas preguntas harían falta para que la semianchura del intervalo fuera de 1 punto porcentual?
\end{enumerate}
\fuente{apartado (b) adaptado de MIT 18.05, primavera de 2014, examen~2, problema~3~\cite{mit1805ex2} (tabla de Bayes discreta para el tipo de dado, con otros dados y otra pregunta predictiva); apartados (a) y (c), de elaboración propia.}
\end{ejercicio}

% ---------------------------------------------------------------------
\begin{ejercicio}{2}{eje:p6}
\textbf{Entropía, compresión y bits por byte (M06). 11 puntos.}

Una fuente emite los símbolos \texttt{a}, \texttt{b}, \texttt{c}, \texttt{d}, \texttt{e} con probabilidades $p=(1/4,\,1/4,\,1/4,\,1/8,\,1/8)$.
\begin{enumerate}[label=(\alph*)]
  \item \pts{3} Calcule $H(p)$ en bits. Construya un código de Huffman, calcule su longitud media y compárela con $H(p)$. Compruebe la desigualdad de Kraft. Codifique el mensaje «\texttt{abed}» e indique cuántos bits ocupa.
  \item \pts{3} Sean $q$ la distribución uniforme sobre los cinco símbolos y $q'$ la distribución $(1/8,\,1/8,\,1/4,\,1/4,\,1/4)$. Calcule $H(p,q)$, $\KL{p}{q}$, $H(p,q')$ y $\KL{p}{q'}$. ¿Cuántos bits por símbolo de más paga cada modelo? ¿Cuál es mejor?
  \item \pts{5} \emph{Estimación.} Un LLM obtiene una pérdida de validación de $1{,}8$ nats por token, con $4{,}2$ bytes de texto por token en promedio. Calcule su perplejidad, los bits por token y los bits por byte. ¿Cuánto ocuparía 1~GB ($10^{9}$ bytes) de ese texto comprimido mediante codificación aritmética con este modelo? ¿Qué razón de compresión supone?
\end{enumerate}
\fuente{apartado (a) adaptado de la hoja de ejemplos de \emph{Information Theory}, Universidad de Cambridge, curso 2025--26, supervisión~1, ejercicio~12~\cite{camit} (códigos de Huffman y comparación de la longitud media con la entropía), según el planteamiento de Cover y Thomas~\cite[§5.6]{ct}; apartados (b) y (c), de elaboración propia.}
\end{ejercicio}
